\documentclass[12pt, a4paper]{article}

% --- Paquetes de Idioma y Formato ---
\usepackage[utf8]{inputenc}
\usepackage[spanish]{babel}
\usepackage[margin=2.5cm]{geometry}
\usepackage{graphicx}
\usepackage{enumitem}
\usepackage{hyperref}

% --- Configuración de Párrafos ---
\setlength{\parskip}{1em}
\setlength{\parindent}{0pt}

% --- Inicio del Documento ---
\begin{document}
\begin{titlepage}
    \centering
    \vspace*{2cm}
    
    % Título Principal
    {\Huge \textbf{Hito 1: Sistemas Avanzados de Integración de la Información}\par}
    
    \vspace{1cm} % Aumentamos el espacio
    
    % Subtítulo más grande (LARGE) y con itálica para dar estilo
    {\LARGE \textit{Sistema de integración de datos climáticos históricos y actuales en España}\par}
    
    \vfill    
    \vfill
    
    % Bloque del Grupo (con letras más grandes)
    {\Large \textbf{Grupo 5}} \\
    \vspace{0.2cm}
    \rule{0.3\linewidth}{0.5pt} % Una línea sutil para destacar el grupo
    \vspace{0.5cm}
    
    \textbf{Autores:}\\
    {\large
    Julio Herrera Ruiz \\
    Irene Retortillo González \\
    Tomás Ruiz Rojo \\
    }
    
    \vspace{1cm}
    
    \textbf{Fecha:}\\
    {\large \today}
    
    \vfill
\end{titlepage}
% --- Índice ---
\tableofcontents
\newpage

\section{Tema y Supuesto elegido: Información sobre el clima}
El supuesto elegido fue el de información sobre el clima, teniendo como tema principal la integración de datos climáticos históricos y actuales en España.
El objetivo de este proyecto es proporcionar información climática, tanto histórica como actual, para España con el objetivo de que se pueda contrastar junto con condiciones geográficas y demográficas,
proporcionando herramientas útiles para la población en general para que puedan tomar decisiones informadas sobre sus actividades diarias, viajes, etc.

Las consultas que hemos elegido realizar para este supuesto son:
\begin{itemize}
    \item Datos de radiación ultravioleta y calidad del aire en provincias de Cataluña.
    \item Datos de temperatura histórica media en marzo en estaciones de la provincia de Asturias en comparación con la predicción para el día siguiente y en relación con la altura del municipio en que se encuentra la estación.
    \item Última predicción de temperatura de la AEMET y de eltiempo.es para municipios de más de 5.000 habitantes en la provincia de Valladolid.
\end{itemize}

\section{Esquema del mediador}














\end{document}